% GENERATED FILE. Edit canonical lessons, then run npm run generate:books.
% canonical-source-hash: fnv1a64:441420c5b6ee87a4
% canonical-lessons: TA-C471-panpu, TA-C471-cupavam, TA-C471-iyalpu, TA-C471-pakkuvam, TA-C471-alatciyam

\chapter{Character and Nature}
\label{ch:ta-character-and-nature}
\hlchaptermodality{471}

\begin{chapteropening}
\textbf{By the end of this chapter:} \emph{I can talk about virtues, temperament, maturity and carelessness.}

Complete the last lesson of chapter 471: I can talk about virtues, temperament, maturity and carelessness.
\end{chapteropening}

\section[\texorpdfstring{\ta{பண்பு} (paṇpu)}{பண்பு (paṇpu)}]{\ta{பண்பு} (paṇpu) — a virtue; good character}

\label{lesson:TA-C471-panpu}



\begin{quote}
Before the new one: say the Tamil for forgiveness, then the Tamil for the cold shoulder.
\end{quote}

\subsection*{You'll want to know: \ta{பண்பு}}
\textbf{\ta{பண்பு}} — \emph{paṇpu} — \textquotedblleft{}a virtue; good character\textquotedblright{}.

\begin{quote}
\ta{உண்மை} \ta{பேசுவது} \ta{ஒரு} \ta{நல்ல} \ta{பண்பு}.
\end{quote}

\begin{quote}
\emph{uṇmai pēcuvatu oru nalla paṇpu.}
\end{quote}

\begin{quote}
Speaking the truth is a good quality.
\end{quote}

\ta{பண்பான}, well-mannered.

\begin{cousinweb}
\textbf{\ta{பண்பு}} (\emph{paṇpu}) is native Dravidian (DEDR 3904), a word for quality, nature and good character.
\end{cousinweb}

\begin{grammarlens}[title={What you've built}]
One more word for describing what you lived and felt.
\end{grammarlens}

\subsection*{Your turn}
\begin{itemize}
\raggedright
  \item \emph{Say it:} \emph{paṇpu}
  \item \emph{Say it:} \emph{paṇpu}, once more
  \item \emph{Say it:} \emph{uṇmai pēcuvatu oru nalla paṇpu.} — Speaking the truth is a good quality.
  \item \emph{Recall:} say the Tamil for forgiveness, then the Tamil for the cold shoulder, then say \emph{paṇpu} again
\end{itemize}

\begin{tcolorbox}[breakable,colback=teal!4,colframe=teal!35!black,title={Before you move on}]
What does \ta{பண்பு} mean? (\textquotedblleft{}A virtue; good character\textquotedblright{}.) Say it once more.
\end{tcolorbox}
\section[\texorpdfstring{\ta{சுபாவம்} (cupāvam)}{சுபாவம் (cupāvam)}]{\ta{சுபாவம்} (cupāvam) — temperament; disposition}

\label{lesson:TA-C471-cupavam}



\begin{quote}
Before the new one: say the Tamil for the cold shoulder, then the Tamil for a virtue.
\end{quote}

\subsection*{You'll want to know: \ta{சுபாவம்}}
\textbf{\ta{சுபாவம்}} — \emph{cupāvam} — \textquotedblleft{}temperament; disposition\textquotedblright{}.

\begin{quote}
\ta{என்} \ta{அண்ணனுக்கு} \ta{அமைதியான} \ta{சுபாவம்}.
\end{quote}

\begin{quote}
\emph{eṉ aṇṇaṉukku amaitiyāṉa cupāvam.}
\end{quote}

\begin{quote}
My elder brother has a calm temperament.
\end{quote}

Dative: \ta{அவனுக்கு} \ta{அமைதியான} \ta{சுபாவம்}, he is calm by nature.

\begin{cousinweb}
\textbf{\ta{சுபாவம்}} (\emph{cupāvam}) is borrowed from Sanskrit \emph{svabhāva}, \textquotedblleft{}one's own nature\textquotedblright{} (Monier-Williams); Tamil drops the \emph{v} of \emph{sv}.
\end{cousinweb}

\begin{grammarlens}[title={What you've built}]
One more word for describing what you lived and felt.
\end{grammarlens}

\subsection*{Your turn}
\begin{itemize}
\raggedright
  \item \emph{Say it:} \emph{cupāvam}
  \item \emph{Say it:} \emph{cupāvam}, once more
  \item \emph{Say it:} \emph{eṉ aṇṇaṉukku amaitiyāṉa cupāvam.} — My elder brother has a calm temperament.
  \item \emph{Recall:} say the Tamil for the cold shoulder, then the Tamil for a virtue, then say \emph{cupāvam} again
\end{itemize}

\begin{tcolorbox}[breakable,colback=teal!4,colframe=teal!35!black,title={Before you move on}]
What does \ta{சுபாவம்} mean? (\textquotedblleft{}Temperament; disposition\textquotedblright{}.) Say it once more.
\end{tcolorbox}
\section[\texorpdfstring{\ta{இயல்பு} (iyalpu)}{இயல்பு (iyalpu)}]{\ta{இயல்பு} (iyalpu) — one’s nature; normality}

\label{lesson:TA-C471-iyalpu}



\begin{quote}
Before the new one: say the Tamil for a virtue, then the Tamil for temperament.
\end{quote}

\subsection*{You'll want to know: \ta{இயல்பு}}
\textbf{\ta{இயல்பு}} — \emph{iyalpu} — \textquotedblleft{}one’s nature; normality\textquotedblright{}.

\begin{quote}
\ta{புதியவர்களிடம்} \ta{வெட்கப்படுவது} \ta{அவளுடைய} \ta{இயல்பு}.
\end{quote}

\begin{quote}
\emph{putiyavarkaḷiṭam veṭkappaṭuvatu avaḷuṭaiya iyalpu.}
\end{quote}

\begin{quote}
Being shy with strangers is her nature.
\end{quote}

\ta{இயல்பாக}, naturally; \ta{இயல்பு} \ta{நிலை}, normal conditions.

\begin{cousinweb}
\textbf{\ta{இயல்பு}} (\emph{iyalpu}) is native Dravidian, a noun from \ta{இயல்} \emph{iyal} (DEDR 471), \textquotedblleft{}to be possible, to befall\textquotedblright{}.
\end{cousinweb}

\begin{grammarlens}[title={What you've built}]
One more word for describing what you lived and felt.
\end{grammarlens}

\subsection*{Your turn}
\begin{itemize}
\raggedright
  \item \emph{Say it:} \emph{iyalpu}
  \item \emph{Say it:} \emph{iyalpu}, once more
  \item \emph{Say it:} \emph{putiyavarkaḷiṭam veṭkappaṭuvatu avaḷuṭaiya iyalpu.} — Being shy with strangers is her nature.
  \item \emph{Recall:} say the Tamil for a virtue, then the Tamil for temperament, then say \emph{iyalpu} again
\end{itemize}

\begin{tcolorbox}[breakable,colback=teal!4,colframe=teal!35!black,title={Before you move on}]
What does \ta{இயல்பு} mean? (\textquotedblleft{}One’s nature; normality\textquotedblright{}.) Say it once more.
\end{tcolorbox}
\section[\texorpdfstring{\ta{பக்குவம்} (pakkuvam)}{பக்குவம் (pakkuvam)}]{\ta{பக்குவம்} (pakkuvam) — maturity; the right stage}

\label{lesson:TA-C471-pakkuvam}



\begin{quote}
Before the new one: say the Tamil for temperament, then the Tamil for one’s nature.
\end{quote}

\subsection*{You'll want to know: \ta{பக்குவம்}}
\textbf{\ta{பக்குவம்}} — \emph{pakkuvam} — \textquotedblleft{}maturity; the right stage\textquotedblright{}.

\begin{quote}
\ta{அவளுக்கு} \ta{வயது} \ta{குறைவு}, \ta{ஆனால்} \ta{நல்ல} \ta{பக்குவம்} \ta{இருக்கிறது}.
\end{quote}

\begin{quote}
\emph{avaḷukku vayatu kuṟaivu, āṉāl nalla pakkuvam irukkiṟatu.}
\end{quote}

\begin{quote}
She is young, but very mature.
\end{quote}

Of food, the right degree of cooking: \ta{சமையல்} \ta{பக்குவம்}.

\begin{cousinweb}
\textbf{\ta{பக்குவம்}} (\emph{pakkuvam}) is borrowed from Sanskrit \emph{pakva}, \textquotedblleft{}cooked, ripe\textquotedblright{} (CDIAL 7621): a ripened mind.
\end{cousinweb}

\begin{grammarlens}[title={What you've built}]
One more word for describing what you lived and felt.
\end{grammarlens}

\subsection*{Your turn}
\begin{itemize}
\raggedright
  \item \emph{Say it:} \emph{pakkuvam}
  \item \emph{Say it:} \emph{pakkuvam}, once more
  \item \emph{Say it:} \emph{avaḷukku vayatu kuṟaivu, āṉāl nalla pakkuvam irukkiṟatu.} — She is young, but very mature.
  \item \emph{Recall:} say the Tamil for temperament, then the Tamil for one’s nature, then say \emph{pakkuvam} again
\end{itemize}

\begin{tcolorbox}[breakable,colback=teal!4,colframe=teal!35!black,title={Before you move on}]
What does \ta{பக்குவம்} mean? (\textquotedblleft{}Maturity; the right stage\textquotedblright{}.) Say it once more.
\end{tcolorbox}
\section[\texorpdfstring{\ta{அலட்சியம்} (alaṭciyam)}{அலட்சியம் (alaṭciyam)}]{\ta{அலட்சியம்} (alaṭciyam) — negligence; indifference}

\label{lesson:TA-C471-alatciyam}



\begin{quote}
Before the new one: say the Tamil for one’s nature, then the Tamil for maturity.
\end{quote}

\subsection*{You'll want to know: \ta{அலட்சியம்}}
\textbf{\ta{அலட்சியம்}} — \emph{alaṭciyam} — \textquotedblleft{}negligence; indifference\textquotedblright{}.

\begin{quote}
\ta{உடல்நலத்தில்} \ta{அலட்சியம்} \ta{கூடாது}.
\end{quote}

\begin{quote}
\emph{uṭalnalattil alaṭciyam kūṭātu.}
\end{quote}

\begin{quote}
One must not be careless about one’s health.
\end{quote}

\ta{அலட்சியமாக}, carelessly.

\begin{cousinweb}
\textbf{\ta{அலட்சியம்}} (\emph{alaṭciyam}) is borrowed from Sanskrit \emph{alakṣya}, \textquotedblleft{}invisible, unnoticed\textquotedblright{} (CDIAL 698); Tamil writes \emph{kṣ} as \ta{ட்ச}.
\end{cousinweb}

\begin{grammarlens}[title={What you've built}]
One more word for describing what you lived and felt.
\end{grammarlens}

\subsection*{Your turn}
\begin{itemize}
\raggedright
  \item \emph{Say it:} \emph{alaṭciyam}
  \item \emph{Say it:} \emph{alaṭciyam}, once more
  \item \emph{Say it:} \emph{uṭalnalattil alaṭciyam kūṭātu.} — One must not be careless about one’s health.
  \item \emph{Recall:} say the Tamil for one’s nature, then the Tamil for maturity, then say \emph{alaṭciyam} again
\end{itemize}

\begin{tcolorbox}[breakable,colback=teal!4,colframe=teal!35!black,title={Before you move on}]
What does \ta{அலட்சியம்} mean? (\textquotedblleft{}Negligence; indifference\textquotedblright{}.) Say it once more.
\end{tcolorbox}
